\documentclass[11pt]{article}
\usepackage{../common/indexus}
\usepackage{../figkit/ixfig}

\begin{document}

\title{The Indexus Manifesto: Data as a Commons}
\author{\ixbyline}
\hypersetup{pdftitle={The Indexus Manifesto: Data as a Commons}, pdfauthor={Rémi Henriot (AI Co-authored)}}
\date{November 2026}

\maketitle

\begin{abstract}
\noindent \textbf{Abstract. }Data is found through indexes. Each new producer makes an index more useful to readers, and each new reader makes it more valuable to producers. Whoever runs that index can decide which records are found and on what terms. We propose that competing producers publish into a shared index. Independent hosts serve it, and applications choose how to present the same data. We seek a way for readers to check a producer's record and for hosts and producers to be paid for what is read.
\vspace{1em}\end{abstract}

\section*{The problem}
The issue is not that data exists in many places. It is that the indexing of data concentrates.

Data is useful once it can be found: a property sale near an address, the weather at a given hour, an offer within a given distance, the events of a given day. Finding requires an index, a structure that answers ``what is here, and how much of it''. Building and maintaining an index of a whole domain is costly, so in practice one or a few parties do it, and everyone else queries theirs.

An index has increasing returns. Each producer that lists in it makes it more complete for readers, and each reader that queries it makes it more valuable for producers. Beyond a certain size it is more useful than any alternative, and both producers and readers converge on it. This is a network effect, and it applies to the index rather than to the data: the same records could be held by many parties, but they are found through one.

The party that operates the dominant index then sets, by construction, what is found, in which order, under which terms of access and at which price for being listed. It also observes every query. Traffic follows the index, not the producers. Competitors who would not share a database with each other each list with the same operator, because that is where the readers are. Applications that want the data must either rebuild an index of their own, at full cost, or depend on the operator's.

Indexus addresses this point: it makes the index itself shared. The index is a protocol that anyone can host and anyone can query, with rules fixed by a public specification, so that the network effect accrues to the shared index rather than to one operator.

\begin{figure}[H]
    \centering
    \ixnew{figures/silos.tikz}
    \caption{Left, each application indexes the same sources on its own, or depends on one operator's index. Right, the sources are indexed once in a shared index that no single party operates, and every application is a view on it.}
    \label{fig:silos}
\end{figure}

\section*{Principles}
These principles describe the network we aim to build. The final status section separates the running index from the proposed verification, payment and governance mechanisms.

\subsection*{1. Data is produced once and read everywhere}
A record should be published once, by the party who produced it, into an index that anyone can read and anyone can host. Reading it should not require copying it. A new application should start from the data that already exists, not from an empty database and a scraping job.

\subsection*{2. Applications are views, not owners}
An application decides what to show, how and to whom. It should not need to own the data to do so. Consider a person comparing yesterday's walk with a weather map. The application reads the walk from the person's collection, asks the weather producer for the area and hour it needs, and draws the two together. Each producer can update its collection independently. Another application can use the same weather without obtaining the person's walk. The \ixpapertext{client}{client paper} describes how the reader chooses compatible cells. The application is decoupled from the production of the data, and can be replaced without losing it.

\subsection*{3. The index is shared, not operated by one party}
Decentralisation addresses that mechanism directly: the index is not held by any single operator. Businesses in competition can publish their offers into the same collections and host part of the same index. They obtain the same network effect as a third-party index, one place where every offer can be found, compared and summed, while the order of results and the meaning of a total are fixed by a public specification, the same for every participant. Ranking, presentation and the customer relationship remain with the applications and businesses that choose them.

\subsection*{4. A protocol, not a platform}
Indexus is the project behind the proposed Distributed Prefix Index Protocol (DPIP). DPIP gives independent implementations a common set of rules for addressing, reading and writing zones. Several independent implementations must be able to join the same network and pass the same conformance tests, as the implementations of an Internet standard do. No implementation, including the one we write, owns the network. Changes go through public proposals.

\subsection*{5. Exact at every scale}
A map shows a city and a street with the same tool. Their totals should account for every record, including deletions. A parent replaces a child's previous contribution when it receives a newer one. A reader needs to know both what a total counts and how fresh it is.

\subsection*{6. Trust is verified, not granted}
In an open network the node that answers may be unfamiliar to the reader. Our aim is for producers to sign their records and for readers to check a total against those signatures.

\subsection*{7. Private stays private}
A person's history can stay on their own device while the application overlays a public forecast. The weather service still sees the area and hour requested. Someone designing the application must decide whether that disclosure is acceptable and what granularity to request. Producer-controlled access to private collections is a distinct goal.

\subsection*{8. Hosting and producing are paid}
Infrastructure costs money to run, and producers incur costs to keep their records current. Base public data should stay free to read. For premium collections we propose dividing the price of a read between the host that served it and the producer. The \ixpapertext{economics}{economics paper} explores settlement through Bitcoin's Lightning Network.

\subsection*{9. Open and not appropriable}
The current reference node and JavaScript SDK carry the MIT licence. A copyleft reference node was proposed in the original outline. We propose a neutral foundation to steward the DPIP specification and its name. Indexus would remain the project's brand. Client libraries stay permissive so that applications built on the index can be anything their authors want.

\subsection*{10. Sovereign and resilient}
Nodes run on ordinary machines, on commercial providers or at home. A host can resume from its disk. Surviving the permanent loss of that disk also requires a copy elsewhere. Public services and organisations can choose where they host their part of the index. We intend to document the uses we support in a public charter.

\subsection*{11. Measured, not promised}
Every claim about scale, latency or exactness is published with the scripts that reproduce it. What is implemented, what is measured and what is only proposed are labelled as such in every paper.

\section*{Where we are}
A holder writes a record and its receipt to disk before replying. A parent can receive the same child summary twice and count it once. The \ixpapertext{properties}{properties paper} reports the tests. The \ixpapertext{history}{history} traces the work from mobile search in 2014 to the two papers of 2024: how can a reader find records incrementally, and how can independent hosts share the index? Verifiable sums, payments and formal standardisation remain goals.

The original four-part outline described the system from the bottom up:
\begin{enumerate}[leftmargin=1.5em]
  \item \textit{Encoding, Spaces and Incremental Search}: one tree for addresses, collections and records, and how a feed or a map reads it.
  \item \textit{A Peer-to-Peer Protocol for Zones}: who holds which zone, how zones move between nodes, and what an acknowledgement guarantees.
  \item \textit{Exact and Verifiable Aggregates}: how every zone keeps an exact sum, how fast parents follow, and how a reader checks a total.
  \item \textit{Paying for Open Data}: how hosts and producers are paid, and how the protocol becomes a standard with many implementations.
\end{enumerate}

The current series keeps this outline as its historical starting point. It separates the \ixpapertext{whitepaper}{whitepaper}, \ixpapertext{client}{client} and \ixpapertext{network}{network}, followed by \ixpapertext{economics}{economics}, \ixpapertext{properties}{properties and applications}, and \ixpapertext{history}{development history}.

\section*{An invitation}
We are looking for producers who want their data read without giving it away, for hosts who want to run part of a common infrastructure, for developers who want to build applications on data they do not have to collect, and for people who want to write a second implementation and prove that the specification stands on its own. The goal is not one more database. It is to make the next application cheaper to build than the last one, because the data it needs is already there, shared, exact and paid for fairly.

\Needspace{3\baselineskip}
\paragraph{Contributions.} The aims in this manifesto are Rémi Henriot's. AI-assisted editing helped distinguish the principles from the features already implemented. The \ixsection{history}{ai}{project history} dates the contributions of the tools used.

\end{document}
